\newpage
\thispagestyle{fancy}
\phantomsection
\addcontentsline{toc}{chapter}{Conclusion générale}
\mtcaddchapter
\markboth{Conclusion générale}{Conclusion générale}

% Titre dans un cadre centré, sur la même page
\begin{center}
\vspace*{0.5cm}
\begin{tikzpicture}
  \node[
    draw        = primary,
    line width  = 0.8pt,
    fill        = white,
    inner xsep  = 2.0cm,
    inner ysep  = 0.7cm,
    text width  = 10cm,
    align       = center
  ]{
    {\fontsize{16}{20}\selectfont\bfseries\scshape\rmfamily\color{primary}
    Conclusion générale}
  };
\end{tikzpicture}
\vspace{0.8cm}
\end{center}

\justifying
À travers ce projet de fin d'études réalisé au sein de \textbf{SolidWall Consulting}, nous avons conçu et développé \textbf{SolidMaint}, une plateforme web de gestion des contrats de maintenance visant à centraliser et automatiser l'ensemble des processus métier de l'entreprise.

\medskip

Ce travail nous a conduits à parcourir toutes les étapes d'un projet logiciel complet : de l'analyse de l'existant et la spécification des besoins, à la modélisation UML, jusqu'à la réalisation et la validation par les tests. Structuré en quatre sprints selon la méthodologie \textbf{Scrum}, le développement a permis de couvrir progressivement l'ensemble du cycle de vie d'une demande de maintenance, depuis sa soumission jusqu'à sa clôture.

\medskip

Sur le plan \textbf{académique}, ce projet nous a permis d'atteindre une maturité supplémentaire en abordant des problématiques d'une complexité inédite, notamment la conception d'une architecture modulaire, l'intégration de l'intelligence artificielle et la mise en place d'une stratégie de tests complète.

\medskip

Sur le plan \textbf{personnel}, cette expérience a représenté une étape importante dans notre parcours en nous confrontant aux réalités du monde professionnel : adapter nos décisions aux contraintes du terrain, communiquer efficacement avec les différentes parties prenantes et assumer la responsabilité d'un produit destiné à un usage réel.

\medskip

En perspective, plusieurs axes d'amélioration pourraient être envisagés, notamment le déploiement sur une infrastructure cloud pour garantir une haute disponibilité, le développement d'une application mobile pour une accessibilité élargie, ainsi que l'enrichissement des capacités d'intelligence artificielle vers une analyse prédictive des incidents récurrents.